% ---- pg-031 （书页 19） ----

首先,对每只熊,我们描出体重对年龄的图然后用直线连接,这样能够近似给出生长曲线图.\hyperref[fig:1.13]{图 1.13} 是所有 7 只熊的生长曲线图.很明显的是,第 5 只熊相对于其他来说有很大不同.这是外族还是自然变异的?在野外,熊的体重是用磅来度量的.进一步的观测表明,这种情况是由于助手在建立电子数据库时,没有将单位换算成公斤,正确的体重为$(45,66,84,112)$公斤.

\begin{mvfigure}
\mvinc{pg-031}{0.343,0.202,0.704,0.412}
\phantomsection\label{fig:1.13}\mvcap{图 1.13\quad 7 只母灰熊体重的联合生长曲线列}
\end{mvfigure}

因为在一个联合生长曲线图上观测个体变化很难,于是可以逐一画出一系列单个生长曲线,这样相似性和差别都能够很容易地观测出来.\hyperref[fig:1.14]{图 1.14} 给出了 7 条体重生长曲线,有些体重生长曲线是线性的,有些是二次的.

\begin{mvfigure}
\mvinc{pg-031}{0.182,0.514,0.864,0.788}
\phantomsection\label{fig:1.14}\mvcap{图 1.14\quad 单个母灰熊体重的生长曲线列}
\end{mvfigure}

\hyperref[fig:1.15]{图 1.15} 给出了单个母灰熊身长的生长曲线列.有一只熊好像在第 2 年到第 3 年变得越来越短了,但是研究人员知道,测量身长用的卷尺会由于打了镇静剂的熊的姿势脱落,进而导致结果的偏差.

现在我们转向多元数据在二维空间中的两个常见的图示表示法: 星形和切尔诺夫脸.
